\section{First-pass problems}
\label{sec:problems}

But\ldots some features the library was missing:

\begin{enumerate}\setcounter{enumi}{5}
  \item No bugs \quad$\times$
  \item Informative error messages \quad$\times$
  \item Rich tactic support \quad$\times$
  \item Good performance (strange but important in some cases) \quad$\times$
  \item Documentation of all points thus far \quad(Some)
\end{enumerate}

\note{These continue the numbering of the wishlist in \Cref{sec:wishlist}, which
is the talk's device for showing that the first five criteria were met and the
development still went badly. Keep the numbering; it is what makes
\Cref{sec:results}'s rundown legible.}

\textbf{Tactic surface in ITrees up 178\%}. Needed 20 custom tactics (vs 14 from
paco).

\note{\textbf{The tactic counts do not match any current measure and need
replacing.} Measured, the count of distinct coinduction-driving tactic names
goes paco $\to$ old port: $16 \to 21$ in \code{theories/}, and $20 \to 23$ across
the whole development. So the first port needed \emph{23} against paco's
\emph{20}, not 20 against 14. The $178\%$ is correct and is the whole-development
character figure. My guess is that ``20 to 14'' was paco against the
\emph{new} port at an earlier measurement, which is now $20 \to 12$; that belongs
in \Cref{sec:results}, not here.}

Why?
\begin{enumerate}
  \item \code{rocq-coinduction} tactics \textbf{not expressive enough}
  \item Some \code{rocq-coinduction} tactics were \textbf{buggy}
\end{enumerate}
Unusable in certain cases, which required extensive workarounds. Error messages
could get very confusing.

\subsection{Confusing errors (\& bugs)}

\begin{gapbox}
The two error figures from slides 75 and 76. Both are reproducible and
released: see \code{paper-data/old-library-errors/}. The second is the stronger
one, and it is the bug:
\begin{lstlisting}
forall (t1 : itree E R1) (t2 : itree E R2),
  gfp (eqit_mon b1 b2) R1 R2 RR t1 t2 ->
  gfp (eqit_mon b1 b2) R1 R2 RR' t1 t2
\end{lstlisting}
\code{coinduction R CIH} on this goal reports
\code{No such section variable or assumption: R.} The message names \code{R},
the candidate the user just wrote, so it reads as though the user's script is at
fault rather than the goal shape being unsupported.
\end{gapbox}

\subsection{Why debugging was hard}

Debugging my own code was \textbf{hard}. Fixing the bugs was \textbf{also hard}.

Why? \code{rocq-coinduction} uses a reifying OCaml plugin with a complex Rocq
invariant.

\gap{The two quotes from slide 79, which are the sharpest evidence in the talk
that the plugin was the obstacle: \emph{``the key invariant is that [e] should be
convertible to [pT a c r x]''} and \emph{``the OCaml type of [g] is a bit
complicated: [bool -> int -> (int -> constr) -> constr]''}.}
